\chapter{Tehnologii Utilizate}
\label{ch:tehnologii}

Dezvoltarea proiectului \textit{Farm Defense: Harvest Wars} a necesitat o selecție riguroasă a tehnologiilor, având ca obiectiv performanța ridicată, portabilitatea și un ciclu de dezvoltare eficient. Această secțiune detaliază ecosistemul software utilizat, argumentând alegerile tehnice făcute.

\section{Motorul de joc Godot 4.x}
\label{sec:godot}

Pentru partea de client (Game Client), a fost ales motorul de joc Godot Engine, versiunea 4.x \cite{godot_docs}. Godot este un motor \textit{open-source} care oferă o arhitectură bazată pe noduri și scene, fiind extrem de flexibil pentru jocuri 2D și 3D. 

Principalele motive pentru alegerea acestui motor sunt:
\begin{itemize}
    \item \textbf{Suport C\# nativ:} Permite utilizarea unui limbaj de programare modern și performant, comun cu cel utilizat pe server.
    \item \textbf{Sistem de animație procedurală:} Facilitate esențială pentru mișcarea fluidă a unităților.
    \item \textbf{Performanță și mărime redusă:} Motorul este optimizat pentru execuție rapidă și are o amprentă mică pe disc, facilitând distribuția.
\end{itemize}

\section{Ecosistemul .NET și Limbajul C\#}
\label{sec:dotnet}

Întregul sistem (atât clientul, cât și serverul) este dezvoltat folosind limbajul C\# (versiunea 10 și ulterioare) pe platforma .NET \cite{csharp_nutshell}. Utilizarea unui limbaj unic pe ambele părți ale aplicației a permis partajarea modelelor de date (Shared Models) și a redus timpul necesar context-switch-ului între limbaje diferite.

Serverul este implementat folosind ASP.NET Core \cite{aspnet_docs}, beneficiind de performanțele ridicate ale acestui framework în gestionarea cererilor HTTP și a proceselor de fundal necesare matchmaking-ului.

\section{Gestiunea Datelor: SQLite și Entity Framework Core}
\label{sec:sqlite_ef}

Pentru persistența datelor la nivel de server, s-a optat pentru SQLite \cite{sqlite_docs}, un motor de baze de date relaționale \textit{serverless}, integrat direct în aplicație. SQLite este ideal pentru acest proiect datorită simplității în configurare și portabilității bazei de date sub formă de fișier unic.

Interacțiunea cu baza de date este realizată prin intermediul \textit{Object-Relational Mapper}-ului (ORM) Entity Framework Core \cite{ef_core_action}. S-a utilizat abordarea \textit{Code-First}, care permite definirea structurii bazei de date direct prin clase C\# și gestionarea evoluției schemei prin migrații controlate.

\section{Mediul de Dezvoltare și Operare: Ubuntu și Linux}
\label{sec:os}

Atât mediul de dezvoltare, cât și cel de operare pentru server, sunt bazate pe sistemul de operare Ubuntu Linux. Stabilitatea și securitatea oferite de nucleul Linux, împreună cu performanța ridicată a runtime-ului .NET pe această platformă, au asigurat un mediu de lucru robust pentru întregul ciclu de viață al proiectului.
